\documentclass[11pt]{article}
\usepackage[margin=1in]{geometry}
\usepackage{amsmath,amssymb,booktabs,graphicx,float,caption,threeparttable,siunitx,url,hyperref,xcolor}
\sisetup{table-number-alignment=center,detect-weight=true,group-separator={,},group-minimum-digits=4,group-digits=integer}
\captionsetup{font=small,labelfont=bf}
\hypersetup{colorlinks=true,linkcolor=blue!45!black,urlcolor=blue!45!black}
\newcommand{\code}[1]{\texttt{\small #1}}
\emergencystretch=3em
\sloppy
\title{\textbf{Target-library adaptation improves selection,\\but does not preserve a transferable selector}\\[0.35em]
\large Updated computational research report: PE-RankFormer, E25--E32}
\author{}
\date{10 September 2026}
\begin{document}
\maketitle
\begin{abstract}
\noindent
Target-library adaptation enables standalone PE-RankFormer to outperform released
OptiPrime on held-out loci within the external Kim library. Native ordinal fine-tuning
(P) achieves mean efficiency at rank 1 of $0.04178$ against $0.04020$, a
$0.158$-percentage-point improvement and $15.7\%$ reduction in selection regret, with
positive per-seed clustered intervals across three seeds. This is an adaptation result
with unequal target-label exposure, not proof of architectural superiority or independent
cross-library confirmation. E29--E32 revise three interpretations of the original report.
First, the actual initializing checkpoint scores $0.03919$, not the published ensemble's
$0.03952$; the 200-group experiment provides no evidence of harm against its matched start.
Second, three-seed pairwise-loss replication improves validation selection over the tested
utility loss, but an advantage over P or the shared-score control is not established.
Third, preserving the original prediction head does not preserve the deployed selector:
source selection efficiency averages $0.12776$ for P, $0.12607$ for utility S, and $0.11460$
for pairwise S, against $0.13350$ at initialization. The shared-versus-dual-head control also
has a substantial score-scale mismatch. The next research question is therefore whether
label-efficient adaptation can improve the deployed selection rule while preserving
transfer across independently evaluated domains.
\end{abstract}
\tableofcontents
\bigskip

\section{Scope and evaluation populations}
The original source benchmark, external zero-shot panel, and within-library adaptation
test answer different questions. They must not be pooled or described as independent
datasets merely because their partition names differ. All work is computational; no
OptiPrime scores, parameters, features or teachers enter PE-RankFormer training.

\begin{table}[htbp]\centering\small
\caption{Prediction and selection answer different questions. Efficiencies are fractions.
Source selection uses 1,463 informative groups; external zero-shot selection uses 30,475
eligible groups. Adaptation test uses 5,557 groups from the same external library.}
\begin{tabular}{lrrrr}\toprule
& \multicolumn{2}{c}{Pooled Spearman} & \multicolumn{2}{c}{Selection @1}\\
Population / our model & OptiPrime & Ours & OptiPrime & Ours\\\midrule
Source fold 0 / published final ensemble & 0.8690 & 0.9079 & 0.1272 & 0.1364\\
External zero-shot / ordinal-S4D ensemble & 0.6931 & 0.7091 & 0.0369 & 0.0363\\
External adaptation test / P, 3 seeds & 0.6922 & 0.7771 & 0.04020 & 0.04178\\
\bottomrule\end{tabular}
\end{table}

The source final ensemble and the external ordinal-S4D ensemble are different model
configurations. The latter contains five source-fold checkpoints; adaptation starts from
one predetermined checkpoint, \code{r4p2\_ordSSM\_cv1}. P uses the native ordinal objective,
not a newly tested conditional-mean regression loss. S adds a separate deployed selector
trained with a group loss; M trains prediction and selection jointly; shared uses one
prediction-derived score for both objectives; Z freezes the encoder and fits a selector.

The external panel contains 118,187 measurement rows selected from the larger Kim library
after overlap exclusions. Its original zero-target-label comparison remains valid as a
separate result: pooled correlation favors PE-RankFormer, while top-choice efficiency
favors OptiPrime. The library was subsequently used for extensive exploration and
adaptation. E29 test re-evaluations are retrospective audits. The two additional E30
pairwise seeds were evaluated on validation and source fold 0, not target test.

\section{Partition, labels and primary endpoint}
% generated by explore_v2/make_adapt_tables.py -- do not edit by hand
\begin{table}[htbp]\centering\begin{threeparttable}\footnotesize
\caption{The adaptation partition. Whole locus components are assigned, because 13,803 alleles
are reachable by more than one protospacer and 9,893 protospacers install more than one allele,
so decision groups are not independent units.}
\label{tab:partition}
\begin{tabular}{l S[table-format=5.0] S[table-format=5.0] S[table-format=5.0]
                  S[table-format=4.0] S[table-format=5.0] S[table-format=2.2] S[table-format=1.4]}
\toprule
Split & {Components} & {Groups} & {Informative} & {Depth $\geq$5} & {Candidates} &
{Mean RTT} & {Mean eff.} \\
\midrule
train & 10,370 & 19,357 & 15,337 & 4,658 & 72,010 & 24.45 & 0.0233 \\
val & 4,952 & 5,561 & 4,656 & 2,056 & 23,120 & 24.64 & 0.0258 \\
test & 4,954 & 5,557 & 4,675 & 2,061 & 23,044 & 24.55 & 0.0253 \\
\bottomrule
\end{tabular}
\begin{tablenotes}[flushleft]\footnotesize
\item 20,276 components over 30,475 alleles and 30,302 protospacers; 76\% are singletons and the
largest holds 30 groups. Assignment balances group count and depth-$\geq$5 count within
strata of edit class, depth and RTT band -- RTT because it is the geometry axis on which the
models are known to be biased, where an accidental imbalance would confound everything.
\end{tablenotes}\end{threeparttable}\end{table}

E25 assigns whole connected components of the allele--protospacer graph. No allele or
protospacer crosses these splits under the recorded keys. This protects against these
specific overlaps; it is not proof that all homologous sequences or shared assay effects
are independent. Training uses 19,357 decision groups and 72,010 distinct candidates;
validation uses 5,561 groups and 23,120 candidates. The original source-trained checkpoint
and frozen comparator provide the same initialization/reference for all relevant arms.

The primary endpoint is group-mean measured efficiency of the top nominated distinct
design, with a fixed tie rule. Repeated measurements of one design are averaged before
ranking. Regret is the group's best measured efficiency minus the nominated efficiency.
Confidence intervals resample locus components, not individual candidates.
The provisional E25 target was a 10\% reduction of released OptiPrime's regret.

\section{Adaptation beats released OptiPrime on this library}
% generated by explore_v2/make_adapt_tables.py -- do not edit by hand
\begin{table}[htbp]\centering\begin{threeparttable}\footnotesize
\caption{The primary endpoint: achieved efficiency at rank 1 on the held-out test components, 5,557 groups over 4,954 locus clusters. Historical E27 results; this library had prior exploratory use and is not independent confirmation.}
\label{tab:test}\scriptsize
\begin{tabular}{l S[table-format=1] S[table-format=1.5] S[table-format=1.5]
                  S[table-format=+1.5] c S[table-format=+2.1] S[table-format=1.4]}
\toprule
Arm & {Seeds} & {Achieved @1} & {SD} & {$\Delta$ vs.\ OP} & {Seeds $+$} &
{Regret $\downarrow$ (\%)} & {$\rho$} \\
\midrule
S, pairwise & 1 & 0.04205 & {--} & +0.00185 & yes & +18.4 & 0.7107 \\
shared & 3 & 0.04193 & 0.00001 & +0.00173 & yes & +17.1 & 0.7624 \\
S, listnet & 1 & 0.04189 & {--} & +0.00169 & yes & +16.8 & 0.7164 \\
P & 3 & 0.04178 & 0.00004 & +0.00158 & yes & +15.7 & 0.7771 \\
S & 3 & 0.04176 & 0.00010 & +0.00156 & yes & +15.4 & 0.7224 \\
M & 3 & 0.04173 & 0.00026 & +0.00153 & yes & +15.2 & 0.7492 \\
Sgeom & 1 & 0.04156 & {--} & +0.00136 & yes & +13.5 & 0.7183 \\
Z & 1 & 0.04112 & {--} & +0.00092 & yes & +9.1 & 0.6981 \\
G & 1 & 0.03968 & {--} & -0.00052 & \textbf{no} & -5.2 & 0.6052 \\
\midrule
OptiPrime (released) & {--} & 0.04020 & {--} & {--} & {--} & {--} & 0.6922 \\
PE-RankFormer, unadapted & {--} & 0.03952 & {--} & -0.00068 & {--} & -6.8 & 0.7102 \\
random choice & {--} & 0.02323 & {--} & {--} & {--} & {--} & {--} \\
perfect chooser & {--} & 0.05028 & {--} & {--} & {--} & {--} & {--} \\
\bottomrule
\end{tabular}
\begin{tablenotes}[flushleft]\footnotesize
\item Full-budget sequence-adapted arms have per-seed locus-clustered intervals
excluding zero at $p=0.001$; G does not ($p=0.134$), and Z has $p=0.002$. ``Regret
$\downarrow$'' is the reduction of OptiPrime's regret; the plan's provisional target
was 10\%.
\item $\rho$ uses the original prediction head, while achieved @1 uses the deployed
selector. These differ for S, M and Z. Means over seeds are not score ensembles.
The unadapted reference is the published five-checkpoint ordinal-S4D ensemble,
not the single checkpoint used to initialize adaptation. Pairwise has three
validation seeds after E30, but only one historical test seed in this table.
\end{tablenotes}\end{threeparttable}\end{table}

P achieves $0.04178$ versus OptiPrime's $0.04020$: approximately $3.9\%$ higher achieved
efficiency and $15.7\%$ lower regret. Each of its three seed-specific paired intervals
excludes zero. These are not a single confidence interval for an ensemble or for the
entire model-development process. Shared, S and M also exceed the released comparator;
the geometry-only control does not, and the frozen-encoder selector has a smaller gain.

A random choice achieves $0.02323$, and the best measured candidate achieves $0.05028$.
Beating OptiPrime therefore does not mean the selection problem is solved. The test result
also does not establish that the architecture is superior under equal target-label exposure.

\paragraph{The matched starting checkpoint matters.}
E29 measures the exact initializing checkpoint at $0.03919$ on test, with Spearman $0.6981$.
The published unadapted ensemble scores $0.03952$ and $0.7102$. Both comparisons may be
reported, but only the former identifies the change from the adapted model's initialization.
P's adapted Spearman is $0.7771$; the difference from $0.7102$ is an ensemble-to-single-model
comparison, not a matched before/after comparison.

\section{Loss replication and what the controls establish}
% generated by explore_v2/make_adapt_tables.py -- do not edit by hand
\begin{table}[htbp]\centering\begin{threeparttable}\small
\caption{Completed three-seed validation comparison, using selected-epoch training-log values.
The additional pairwise seeds were not evaluated on target test.}
\label{tab:replication}
\begin{tabular}{lrr}\toprule Arm & Seeds & Validation achieved @1 \\
\midrule
P & 3 & 0.042713 \\
S & 3 & 0.042673 \\
S, pairwise & 3 & 0.042998 \\
shared & 3 & 0.043040 \\
\bottomrule\end{tabular}
\end{threeparttable}\end{table}

% generated by explore_v2/make_adapt_tables.py -- do not edit by hand
\begin{table}[htbp]\centering\begin{threeparttable}\small
\caption{Direct contrasts from cached validation predictions, averaging each group's
outcomes over three seeds before resampling whole locus components.}
\label{tab:followupcontrasts}
\begin{tabular}{lrrr}\toprule Contrast & Difference & 95\% interval & $p$ \\
\midrule
Pairwise $-$ P & +0.000291 & $[-0.000072, +0.000654]$ & 0.126 \\
Pairwise $-$ S & +0.000331 & $[+0.000082, +0.000594]$ & 0.009 \\
Pairwise $-$ shared & -0.000023 & $[-0.000270, +0.000233]$ & 0.888 \\
\bottomrule\end{tabular}
\begin{tablenotes}[flushleft]\footnotesize
\item Conditional on selected checkpoints; unadjusted development comparisons, not
independent confirmation. Re-evaluation values can differ slightly from training logs
because of numerical/batch effects. Seed means do not represent an ensemble.
\end{tablenotes}\end{threeparttable}\end{table}

E30 adds two original-recipe pairwise-loss runs, producing three seeds in total.
Selected-epoch validation achieved efficiency averages $0.042998$ for pairwise S,
$0.042713$ for P, $0.042673$ for utility S, and $0.043040$ for shared.
Pairwise exceeds P and utility S in each matched validation seed, but not shared.

Direct cached-prediction contrasts show a pairwise-minus-P difference of $+0.000291$
with interval $[-0.000072,+0.000654]$, $p=0.126$; pairwise-minus-utility-S is
$+0.000331$ with interval $[+0.000082,+0.000594]$, $p=0.009$; and pairwise-minus-shared
is $-0.000023$ with interval $[-0.000270,+0.000233]$, $p=0.888$.
These are unadjusted development comparisons conditional on selected checkpoints.
They support further loss research, not a confirmed overall winner. Numerical/batch
differences between training-log validation maxima and cached re-evaluations are retained
rather than silently forcing them to match.

On the old test split, three-seed utility-S minus P is $-0.000021$ with interval
$[-0.000436,+0.000378]$, $p=0.939$. Nonsignificance does not establish equivalence.
Pairwise's higher historical test point estimate still has only one tested seed.

\paragraph{Head separation was not cleanly isolated.}
Shared adds 199,040 parameters and M adds 198,657, but their architectures and selection-score
ranges also differ. At training temperature 1, the fitted M head has median within-group
score range $87.5$ and mean maximum softmax probability $0.9823$; shared has range $0.2451$
and mean maximum probability $0.3409$. For eight candidates with scores in $[0,1]$,
the maximum possible probability at this temperature is $e/(e+7)=0.2797$.
Thus parameter-count matching does not make the training surrogates comparable.
Post-hoc positive temperature scaling cannot change deterministic argmax decisions;
a repaired comparison must alter training and control capacity, score scale and initialization.

\paragraph{Geometry and candidate-set models remain open hypotheses.}
One scalar-geometry run scored below utility S. This does not reject segment-specific pooling,
validated edit-to-RTT alignment features, or geometry-aware architectures. Candidate-set
attention was not implemented, so it has not been experimentally rejected. Both directions
are lower priority than a well-controlled selection/readout study and independent validation.

\section{Retention of prediction is not retention of selection}
% generated by explore_v2/make_adapt_tables.py -- do not edit by hand
\begin{table}[htbp]\centering\begin{threeparttable}\small
\caption{Both outputs on source fold 0: 20,509 rows and 1,463 informative groups.
Means across available seeds, not ensembles; M, shared and Z are single-seed audits.}
\label{tab:heads}
\begin{tabular}{lr rrrr}\toprule
Arm & Seeds & Prediction @1 & Selection @1 & Prediction $\rho$ & Selection $\rho$ \\
\midrule
Starting checkpoint & 1 & 0.13350 & 0.13350 & 0.8984 & 0.8984 \\
P & 3 & 0.12776 & 0.12776 & 0.8697 & 0.8697 \\
S & 3 & 0.13323 & 0.12607 & 0.8909 & 0.7731 \\
S, pairwise & 3 & 0.13428 & 0.11460 & 0.8917 & 0.7643 \\
M & 1 & 0.12730 & 0.12675 & 0.8768 & 0.7789 \\
shared & 1 & 0.12527 & 0.12527 & 0.8457 & 0.8457 \\
Z & 1 & 0.13350 & 0.12697 & 0.8984 & 0.8110 \\
\bottomrule\end{tabular}
\begin{tablenotes}[flushleft]\footnotesize
\item Source achieved @1 uses informative groups with at least two distinct designs;
target results include all eligible groups. Absolute levels are not directly comparable.
P and shared deploy the prediction score; S, M and Z deploy the separate selection score.
\end{tablenotes}\end{threeparttable}\end{table}

E28 evaluated only the original prediction head. E29 explicitly evaluates both original
prediction and deployed selection outputs on source fold 0. P's source selection averages
$0.12776$; utility S averages $0.12607$, below P in each matched seed, despite retaining a
strong original prediction head. Pairwise S averages $0.11460$ through its deployed selector
while its original prediction head reaches $0.13428$.
These source means are descriptive, not a source superiority or equivalence test.

Z is especially informative: the frozen encoder and original predictor preserve source
@1 at $0.13350$, but the newly fitted selector achieves $0.12697$. Encoder freezing therefore
does not imply that a new deployed selection rule preserves source utility. Readout
specialization is a plausible contributor to the transfer failure; it is not yet a
causally isolated mechanism. The old recommendation to ship S as one selector for many
libraries is withdrawn. Explicit domain-based head routing would be a different deployment
policy and needs its own evaluation.

\section{What the label-budget runs do and do not show}
% generated by explore_v2/make_adapt_tables.py -- do not edit by hand
\begin{table}[htbp]\centering\begin{threeparttable}\footnotesize
\caption{Historical training-budget experiments, on the test components. Subsets take whole
decision groups, not necessarily whole locus components, and are not consistently nested.}
\label{tab:curve}
\begin{tabular}{S[table-format=5.0] S[table-format=1.5] S[table-format=+1.5]
                S[table-format=+1.5] S[table-format=1.3]}
\toprule
{Training groups} & {Achieved @1} & {vs.\ OptiPrime} & {vs.\ ensemble} & {$p$ vs.\ OptiPrime} \\
\midrule
200 & 0.03929 & -0.00091 & -0.00023 & 0.006 \\
1,000 & 0.04021 & +0.00001 & +0.00069 & 0.971 \\
5,000 & 0.04137 & +0.00117 & +0.00185 & 0.001 \\
19,357 & 0.04178 & +0.00158 & +0.00226 & 0.001 \\
\bottomrule
\end{tabular}
\begin{tablenotes}[flushleft]\footnotesize
\item References: OptiPrime 0.04020, published unadapted ensemble 0.03952. Every budget
also used 5,561 labelled validation groups. Small budgets have one seed/subsample;
the full budget has three seeds. These are not total-label learning curves.
The matched-start audit does not establish harm at 200 training groups.
\end{tablenotes}\end{threeparttable}\end{table}

The 200-group run scores $0.03929$, giving $+0.000102$ against the actual starting checkpoint,
with interval $[-0.000041,+0.000255]$, $p=0.158$. It establishes neither harm nor a reliable
benefit. Comparing it with a stronger unadapted ensemble did not isolate fine-tuning harm.

Every historical budget additionally used 5,561 labelled validation groups. The apparent
crossover between 1,000 and 5,000 training groups is specific to those runs; it is not a
total-label threshold. Small budgets had one seed/subsample and unequal numbers of optimizer
updates under an epoch-based schedule. The 200-group subset shares only 12 groups with the
1,000-group subset. Training subsets include partial training components, although the outer
train/validation/test partition remains component-disjoint.

E31 has constructed, but not trained, corrected nested complete-component budgets for three
subset seeds. Nominal 200-group total budgets contain 200--202 groups, with inner validation
allocated inside each selected budget. Candidate measurements, components, training labels
and validation labels are counted separately. Future low-budget runs must allow the
unadapted initialization as an eligible checkpoint and distinguish fixed-update from
fixed-epoch comparisons.

\section{Current claims and implications for the next research phase}
\begin{enumerate}
\item Standalone target adaptation outperforms released OptiPrime on held-out loci within
this library. The source-only external selection deficit remains a separate finding.
\item P remains a strong simple adaptation baseline and shared a strong target-selection
control. Pairwise loss is useful to investigate, but is not established as superior to P
or shared and has a substantial source-selection cost.
\item Source retention must be evaluated through the actual deployed score. Neither an
unused prediction head nor a frozen encoder guarantees this.
\item Architecture, loss, label efficiency and domain transfer require matched controls.
The current data do not support rejecting all specialized losses or dual-head designs.
\item All adapted models saw target labels; released OptiPrime did not. A later
matched-label comparison is distinct from any hybrid and remains necessary before an
architectural-superiority claim. No hybrid is proposed here.
\item One external library in HEK293T/PE2 is insufficient to establish broad transfer.
A stronger methods contribution needs independent computational domains, realistic
label budgets and a reproducible selection protocol, not merely a small gain on this panel.
\end{enumerate}

\section*{Reproducibility and completed work}
E25--E28 contain twenty original training runs. E30 adds two pairwise replications.
E29 provides 42 separately cached model-by-surface prediction tables, with full SHA-256
hashes of recorded checkpoint/data inputs. E31 supplies budget manifests; E32 supplies a
score-scale diagnostic, not new architecture training. Historical experiment JSON and
checkpoints are preserved. Tables are generated by
\path{explore_v2/make_adapt_tables.py} from E25--E32 artifacts.
The current endpoint/protocol suite has 21 passing tests.
Agreement with saved numbers guards transcription, not interpretation or experimental validity.
The detailed current ledger is \path{explore_v2/ADAPTATION_FOLLOWUP_RESULTS.md}.
The proposed next programme is \path{explore_v3/RESEARCH_PLAN.md}; its experiments
have not yet been executed.
\end{document}
